\uuid{PS5z}
\titre{Primitives usuelles}
\auteur{hanani ; mzari}
\organisation{exo7}
\licence{CC BY-NC-SA 4.0}
\datecreate{2020-01-20}
\niveau{L1}
\module{Analyse}
\chapitre{Calcul d'intégrales}
\sousChapitre{Primitives diverses}
\difficulte{2}
\theme{Primitives, primitives usuelles}
\exo7qcmid{lille-2.primitives.015}

\contenu{
\begin{question}
\qclassification{127.11}
\qtags{motcle=primitives usuelles}
Parmi les égalités suivantes, cocher celles qui sont vraies :
\begin{answers}
\bad{$\displaystyle \int \left(\mathrm{e}^{3x}+\frac{1}{1+4x^2}\right)\mathrm{d}x=\mathrm{e}^{3x}+\arctan (2x)+k,\; k\in \Rr$.}
\good{$\displaystyle \int \left(\mathrm{e}^{3x}+\frac{1}{1+4x^2}\right)\mathrm{d}x=\frac{\mathrm{e}^{3x}}{3}+\frac{\arctan (2x)}{2}+k,\; k\in \Rr$.}
\good{$\displaystyle \int \left(\mathrm{e}^{3x}+\frac{4}{1+x^2}\right)\mathrm{d}x=\frac{\mathrm{e}^{3x}}{3}+4\arctan x+k,\; k\in \Rr$.}
\bad{$\displaystyle \int \left(\mathrm{e}^{3x}+\frac{4}{1+x^2}\right)\mathrm{d}x=\mathrm{e}^{3x}+4\arctan x+k,\; k\in \Rr$.}
\end{answers}
\begin{explanations}
Utiliser la linéarité en remarquant que : $\displaystyle \int \mathrm{e}^{3x}\mathrm{d}x=\frac{\mathrm{e}^{3x}}{3}+k_1$,
$$\int \frac{\mathrm{d}x}{1+4x^2}=\frac{\arctan (2x)}{2}+k_2\quad \mbox{et}\quad \int \frac{4\, \mathrm{d}x}{1+x^2}=4\arctan x+k_3,\; k_i\in \Rr.$$
\end{explanations}
\end{question}
}
